\chapter{Cơ sở lý thuyết và tổng quan công nghệ}

Chương này trình bày nền tảng lý thuyết của bài toán cá nhân hóa thời trang, truy hồi tăng cường cho sinh nội dung, thử đồ ảo và kiến trúc xử lý bất đồng bộ. Trên cơ sở đó, các công nghệ được lựa chọn cho hệ thống được tổng hợp theo vai trò kỹ thuật.

\medskip

\section{Cá nhân hóa trải nghiệm trong thương mại điện tử thời trang}

Cá nhân hóa là một yêu cầu quan trọng trong thương mại điện tử thời trang vì quyết định mua hàng không chỉ phụ thuộc vào thuộc tính sản phẩm mà còn chịu ảnh hưởng của sở thích, vóc dáng, hoàn cảnh sử dụng và khả năng kết hợp giữa các trang phục. Các hệ thống gợi ý thời trang vì vậy thường khai thác đồng thời đặc trưng hình ảnh, thuộc tính sản phẩm, lịch sử tương tác và thông tin ngữ cảnh của người dùng \cite{fashion-review}.

\medskip

Trong phạm vi đề tài, dữ liệu cá nhân hóa gồm hồ sơ cơ thể, ảnh người dùng và tủ đồ cá nhân. Những dữ liệu này được kết hợp với danh mục sản phẩm của cửa hàng để hỗ trợ hai nhiệm vụ: lựa chọn phương án phối đồ phù hợp với yêu cầu tự nhiên và tạo ảnh thử đồ ảo.

\medskip

Khác với các hệ thống gợi ý chỉ dựa trên lịch sử nhấp chuột hoặc lịch sử mua hàng, bài toán trong đề tài còn phải xử lý ngữ cảnh thời trang rất giàu nghĩa. Một câu hỏi như ``đi phỏng vấn nên mặc gì'' không chỉ là truy vấn tìm sản phẩm áo sơ mi, mà còn liên quan tới mức độ trang trọng, màu sắc an toàn, kiểu dáng ít gây phân tán và khả năng phối cùng những món đồ người dùng đã có. Vì vậy, hệ thống cần kết hợp dữ liệu có cấu trúc như \code{category}, \code{price}, \code{gender} với dữ liệu ngữ nghĩa như phong cách, chất liệu, hoàn cảnh sử dụng và lời diễn đạt tự nhiên của người dùng.

\medskip

Hệ thống không sử dụng hồ sơ cơ thể để suy luận kích thước trang phục hoặc mô phỏng độ vừa vặn vật lý. Thông tin này được sử dụng như ngữ cảnh cho quá trình tư vấn và làm nguồn ảnh đầu vào cho tác vụ thử đồ ảo. Phạm vi trên giúp tách rõ bài toán gợi ý và sinh ảnh khỏi bài toán dự đoán kích cỡ.

\medskip

Để gợi ý có thể chuyển thành hành động mua sắm, nội dung tư vấn cần được ràng buộc vào tập sản phẩm và trang phục thực sự tồn tại. Cơ chế truy hồi tăng cường cho sinh nội dung đáp ứng yêu cầu này bằng cách cung cấp dữ liệu liên quan cho mô hình trước khi tạo câu trả lời.

\section{AI Stylist và hệ thống RAG}

Retrieval-Augmented Generation (RAG) kết hợp bộ nhớ tham số của mô hình ngôn ngữ với nguồn dữ liệu ngoài được truy hồi tại thời điểm xử lý \cite{rag}. Trong bài toán AI Stylist, bước truy hồi tạo một tập ứng viên từ danh mục sản phẩm và tủ đồ cá nhân; mô hình ngôn ngữ lựa chọn và giải thích phương án phối đồ dựa trên tập ứng viên này. Cách tổ chức đó làm giảm nguy cơ đề xuất sản phẩm không tồn tại và cho phép cập nhật danh mục mà không cần huấn luyện lại mô hình.

\medskip

Luồng xử lý gồm ba bước. Trước hết, yêu cầu tự nhiên được phân tích để xác định nhóm sản phẩm, hoàn cảnh sử dụng, phong cách, khoảng giá và từ khóa tìm kiếm. Tiếp theo, từ khóa được biểu diễn bằng vector nhúng 512 chiều và truy vấn trên cơ sở dữ liệu vector bằng độ tương đồng cosine kết hợp bộ lọc thuộc tính \cite{sentence-transformers,lancedb}. Cuối cùng, mô hình ngôn ngữ nhận tập ứng viên đã truy hồi để sinh các phương án phối đồ có cấu trúc.

\medskip

Đầu ra có cấu trúc cho phép giao diện ánh xạ từng trang phục thành hành động cụ thể như thêm vào giỏ hàng, tìm sản phẩm tương tự hoặc thử đồ ảo. Khác với câu trả lời văn bản thuần túy, cấu trúc JSON còn giúp hệ thống nghiệp vụ kiểm tra định danh và bổ sung dữ liệu giao dịch trước khi trả kết quả cho người dùng.

\medskip

RAG phù hợp với danh mục sản phẩm thay đổi theo thời gian vì dữ liệu mới có thể được đưa vào kho truy hồi độc lập với mô hình ngôn ngữ. Bộ lọc thuộc tính còn cho phép giới hạn kết quả theo danh mục, khoảng giá hoặc trạng thái tồn kho trước khi mô hình tạo phương án phối đồ. Dù vậy, chất lượng đầu ra vẫn phụ thuộc vào độ đầy đủ của siêu dữ liệu, chất lượng vector nhúng và khả năng tuân thủ cấu trúc của mô hình.

\medskip

Trong hệ thống thương mại điện tử, chất lượng của luồng RAG không chỉ thể hiện ở độ tự nhiên của lời tư vấn. Mỗi trang phục được chọn còn phải tồn tại trong nguồn dữ liệu, có hình ảnh và biến thể hợp lệ, đồng thời hỗ trợ các hành động tiếp theo trên giao diện. Vì vậy, dữ liệu do mô hình sinh ra cần được kiểm tra với nguồn dữ liệu giao dịch trước khi hiển thị.

\medskip

Hình~\ref{fig:rag-flow} mô tả quá trình chuyển đổi từ yêu cầu của người dùng thành dữ liệu gợi ý có thể hiển thị trên giao diện.

\begin{figure}[H]
\centering
\resizebox{0.92\textwidth}{!}{%
\begin{tikzpicture}[node distance=11mm and 13mm]
    \node[diagram user, minimum width=3.1cm] (query) {Truy vấn tiếng Việt\\của người dùng};
    \node[diagram ai, right=of query, minimum width=3.4cm] (extract) {\textbf{1. Extract}\\Gemini phân tích intent,\\guardrail và bộ lọc};
    \node[diagram service, right=of extract, minimum width=3.2cm] (embed) {CLIP multilingual\\sinh vector 512 chiều};
    \node[diagram data, right=of embed, minimum width=3.7cm] (retrieve) {\textbf{2. Retrieve}\\LanceDB hybrid search\\\code{products_vector}\\+\code{closet_vector}};

    \node[diagram ai, below=of retrieve, minimum width=3.7cm] (generate) {\textbf{3. Generate}\\Gemini chỉ chọn item trong pool\\và sinh outfit JSON};
    \node[diagram service, left=of generate, minimum width=3.5cm] (hydrate) {Medusa kiểm tra ID,\\hydrate dữ liệu thật\\và lưu session};
    \node[diagram user, left=of hydrate, minimum width=3.3cm] (ui) {Frontend hiển thị\\OutfitCard và hành động};

    \draw[diagram arrow] (query) -- (extract);
    \draw[diagram arrow] (extract) -- (embed);
    \draw[diagram arrow] (embed) -- (retrieve);
    \draw[diagram arrow] (retrieve) -- (generate);
    \draw[diagram arrow] (generate) -- (hydrate);
    \draw[diagram arrow] (hydrate) -- (ui);

    \node[diagram note, above=11mm of extract, minimum width=3.4cm] (refusal) {Nếu ngoài miền thời trang:\\trả lời từ chối, không truy hồi};
    \draw[diagram async] (extract) -- (refusal);

    \node[diagram data, below=11mm of hydrate, minimum width=3.6cm] (source) {Product Module +\\AI Personalization Module};
    \draw[diagram arrow] (source) -- node[right, font=\scriptsize] {title, ảnh, variant, category} (hydrate);
\end{tikzpicture}%
}
\caption{Sơ đồ luồng RAG của AI Stylist theo mã nguồn triển khai}
\label{fig:rag-flow}
\end{figure}


\medskip

Sau khi hệ thống có khả năng đề xuất sản phẩm phù hợp, thử đồ ảo giúp người dùng hình dung trang phục khi mặc lên ảnh của chính mình.

\section{Virtual Try-On trong thời trang số}

Thử đồ ảo (Virtual Try-On -- VTON) là bài toán sinh ảnh người dùng mặc một trang phục tham chiếu. Đầu vào thường gồm ảnh người, ảnh trang phục và một số tham số xử lý như loại trang phục, mặt nạ hoặc độ phân giải. Đầu ra là ảnh mới giữ lại đặc điểm cơ thể, khuôn mặt, dáng đứng của người dùng và thay đổi vùng trang phục theo ảnh tham chiếu \cite{catvton}.

\medskip

CatVTON là mô hình thử đồ ảo dựa trên khuếch tán, sử dụng phép ghép ảnh người và ảnh trang phục theo chiều không gian để đơn giản hóa quá trình suy luận \cite{catvton}. Bên cạnh mô hình chuyên biệt chạy cục bộ, dịch vụ sinh và chỉnh sửa ảnh trên đám mây có thể tiếp nhận nhiều ảnh tham chiếu để tạo kết quả toàn bộ trang phục trong một yêu cầu \cite{replicate}.

\medskip

Hai phương án thể hiện những đánh đổi khác nhau. Mô hình cục bộ cho phép chủ động kiểm soát dữ liệu và tham số suy luận nhưng phụ thuộc tài nguyên phần cứng. Dịch vụ đám mây giảm tải cho máy cục bộ và hỗ trợ nhiều ảnh tham chiếu, đổi lại phụ thuộc kết nối mạng, giới hạn dịch vụ và chi phí sử dụng.

\medskip

VTON xử lý ảnh, vùng mặt nạ, độ phân giải và tài nguyên GPU nên thường có độ trễ lớn hơn yêu cầu web thông thường. Kết quả còn phụ thuộc vào chất lượng ảnh người dùng, góc chụp, ảnh trang phục và mô hình được lựa chọn. Do đó, thử đồ ảo trong đề tài được xem là công cụ hỗ trợ hình dung trang phục, không phải phép đo kích thước hoặc mô phỏng vật lý chính xác.

\medskip

Bảng~\ref{tab:vton-approaches} tổng hợp vai trò và đánh đổi của hai phương án xử lý được lựa chọn.

\begin{table}[H]
\centering
\caption{So sánh hai hướng xử lý Virtual Try-On trong hệ thống}
\label{tab:vton-approaches}
\begin{tabularx}{\textwidth}{L{3.1cm}YY}
\toprule
\textbf{Tiêu chí} & \textbf{CatVTON cục bộ} & \textbf{FLUX.2 đám mây} \\
\midrule
Môi trường xử lý & Chạy trong dịch vụ AI bằng PyTorch và Diffusers & Gửi yêu cầu tới Replicate và nhận địa chỉ ảnh kết quả \\
\addlinespace
Ưu điểm & Chủ động kiểm soát dữ liệu và quá trình suy luận & Giảm tải GPU cục bộ, hỗ trợ nhiều ảnh tham chiếu trong một yêu cầu \\
\addlinespace
Hạn chế & Phụ thuộc VRAM và thời gian suy luận trên máy triển khai & Phụ thuộc kết nối mạng, thông tin xác thực và chính sách dịch vụ \\
\addlinespace
Vai trò & Thử đồ tuần tự theo từng vùng trang phục & Tạo kết quả toàn bộ trang phục từ nhiều ảnh tham chiếu \\
\bottomrule
\end{tabularx}
\end{table}

\medskip

Vì VTON là tác vụ nặng và có độ trễ cao, việc tích hợp nó vào hệ thống thương mại điện tử không chỉ là vấn đề chọn mô hình sinh ảnh. Hệ thống còn cần một kiến trúc xử lý phù hợp để tách các thao tác mua sắm nhanh khỏi các tác vụ AI bất đồng bộ.

\section{Kiến trúc dịch vụ và xử lý bất đồng bộ}

Kiến trúc microservices phân tách hệ thống thành các dịch vụ có trách nhiệm và vòng đời triển khai tương đối độc lập \cite{microservices-study}. Trong đề tài, thao tác thương mại điện tử cần phản hồi nhanh, trong khi sinh ảnh VTON có độ trễ cao và tiêu tốn GPU. Sự khác biệt về tải là cơ sở để tách dịch vụ AI khỏi hệ thống nghiệp vụ thương mại điện tử và sử dụng hàng đợi cho tác vụ sinh ảnh.

\medskip

RabbitMQ đóng vai trò trung gian giữa MedusaJS và dịch vụ AI FastAPI \cite{rabbitmq}. Medusa tạo tác vụ và phát thông điệp vào hàng đợi \code{ai_vision_jobs}. Bộ xử lý AI tiếp nhận thông điệp, xử lý ảnh rồi phát kết quả vào \code{ai_vision_results}. Medusa nhận kết quả, cập nhật trạng thái tác vụ và gửi sự kiện tới giao diện qua SSE.

\medskip

Sự kiện máy chủ gửi tới máy khách (Server-Sent Events -- SSE) được sử dụng vì luồng cập nhật chủ yếu đi một chiều từ máy chủ tới trình duyệt. Cơ chế này dựa trên \code{EventSource} và phản hồi \code{text/event-stream} \cite{sse-mdn}. So với truy vấn định kỳ, SSE giúp giao diện nhận kết quả ngay khi Medusa có dữ liệu mới mà không phải gửi yêu cầu lặp lại.

\medskip

Trong thiết kế bất đồng bộ, hàng đợi không chỉ giúp tránh quá thời gian chờ mà còn tạo ranh giới trách nhiệm giữa các dịch vụ. Medusa tạo tác vụ, lưu trạng thái và phát sự kiện; dịch vụ AI xử lý vector nhúng, siêu dữ liệu và suy luận ảnh. Khi một tác vụ AI bị chậm, giao diện vẫn có thể hiển thị trạng thái đang xử lý. Việc thay đổi bộ xử lý AI cũng ít ảnh hưởng tới phần thương mại điện tử nếu cấu trúc thông điệp và trạng thái tác vụ được giữ ổn định.

\medskip

SSE được chọn thay vì WebSocket vì nhu cầu hiện tại chỉ là máy chủ đẩy kết quả VTON tới trình duyệt. WebSocket phù hợp hơn khi hai phía cần trao đổi liên tục theo thời gian thực. Trong phạm vi bản triển khai tối thiểu, giao diện chỉ cần nhận trạng thái tác vụ và địa chỉ ảnh kết quả nên SSE đơn giản hơn và đủ đáp ứng yêu cầu.

\medskip

Các cơ sở lý thuyết trên được hiện thực bằng tập công nghệ trình bày trong Bảng~\ref{tab:technology-stack}.

\section{Tổng quan công nghệ nền tảng}

\begin{table}[H]
\centering
\small
\caption{Công nghệ nền tảng và vai trò trong hệ thống}
\label{tab:technology-stack}
\begin{tabularx}{\textwidth}{L{2.8cm}L{3.6cm}Y}
\toprule
\textbf{Nhóm công nghệ} & \textbf{Công nghệ} & \textbf{Vai trò} \\
\midrule
Kho mã nguồn hợp nhất & pnpm workspace, Turborepo & Quản lý nhiều ứng dụng và gói dùng chung trong một kho mã nguồn \cite{pnpm,turborepo}. \\
\addlinespace
Giao diện người dùng & Next.js 14, React 18, Tailwind CSS, Zustand & Xây dựng cửa hàng, giao diện AI Stylist, hồ sơ AI, tủ đồ, VTON và giỏ hàng \cite{nextjs,react,tailwind,zustand}. \\
\addlinespace
Hệ thống nghiệp vụ & MedusaJS 2.14, TypeScript & Quản lý sản phẩm, API tùy chỉnh, mô-đun cá nhân hóa và điều phối nghiệp vụ \cite{medusa}. \\
\addlinespace
Cơ sở dữ liệu quan hệ & PostgreSQL & Lưu dữ liệu thương mại điện tử và dữ liệu cá nhân hóa AI \cite{postgresql}. \\
\addlinespace
Hạ tầng & Redis, Docker & Hỗ trợ bộ nhớ đệm và môi trường chạy cục bộ \cite{redis,docker}. \\
\addlinespace
Hàng đợi thông điệp & RabbitMQ & Truyền sự kiện đồng bộ siêu dữ liệu và tác vụ xử lý ảnh giữa các dịch vụ \cite{rabbitmq}. \\
\addlinespace
Dịch vụ AI & FastAPI, Pydantic, Pika & Cung cấp API nội bộ, tiêu thụ thông điệp và thực thi các luồng AI \cite{fastapi}. \\
\addlinespace
Cơ sở dữ liệu vector & LanceDB, PyArrow & Lưu vector nhúng của sản phẩm và tủ đồ cá nhân để tìm kiếm ngữ nghĩa \cite{lancedb}. \\
\addlinespace
Mô hình vector nhúng & \code{clip-ViT-B-32-}\allowbreak\code{multilingual-v1} & Sinh vector 512 chiều cho văn bản và hình ảnh, hỗ trợ truy vấn tiếng Việt \cite{sentence-transformers}. \\
\addlinespace
Mô hình ngôn ngữ & Google Gemini & Phân tích ý định, trích xuất siêu dữ liệu và sinh phương án phối đồ \cite{gemini}. \\
\addlinespace
Thử đồ ảo & CatVTON, FLUX.2 Pro qua Replicate & Sinh ảnh thử đồ bằng mô hình cục bộ hoặc dịch vụ đám mây \cite{catvton,replicate}. \\
\bottomrule
\end{tabularx}
\end{table}

\clearpage
